% ===============================================================
% PEDAGOGIE DES TP IA NUMERIQUE
% Annexe A : choix de conception. Support HTML centralise, formats de
% fichiers, scenarisation, defi et cybersecurite de l'IA
%
% (c) 2026 Selina Cheggour. Tous droits reserves.
% Presentee au conseil de perfectionnement de l'Icam Toulouse,
% 14 et 15 octobre 2026.
%
% Compilation (deux passes) : pdflatex Annexe_A_Choix_de_conception.tex
% ===============================================================
\documentclass[11pt,a4paper]{article}
\usepackage[utf8]{inputenc}
\usepackage[T1]{fontenc}
\usepackage[french]{babel}
\usepackage[default]{opensans}
\usepackage{montserrat}
\usepackage[a4paper,margin=2.3cm,top=2.6cm,bottom=3cm,footskip=1.4cm]{geometry}
\usepackage{xcolor}
\usepackage{tabularx}
\usepackage{booktabs}
\usepackage{array}
\usepackage{enumitem}
\usepackage{fancyhdr}
\usepackage{titlesec}
\usepackage[most]{tcolorbox}
\usepackage{lastpage}

\definecolor{icamorange}{RGB}{243,146,0}
\definecolor{icamred}{RGB}{184,32,24}
\definecolor{icamdark}{RGB}{64,64,64}
\definecolor{icamgray}{RGB}{96,96,96}
\definecolor{fondgris}{RGB}{243,243,241}
\definecolor{fondorange}{RGB}{254,245,232}
\definecolor{encre}{RGB}{43,43,43}

\newcommand{\auteure}{Selina Cheggour}
\newcommand{\droits}{© 2026 \auteure{} · Tous droits réservés · Reproduction, adaptation ou réutilisation interdites sans autorisation écrite de l'auteure}
\newcommand{\presente}{Méthode présentée au conseil de perfectionnement de l'Icam Toulouse, 14 et 15 octobre 2026}

\usepackage[hidelinks]{hyperref}
\hypersetup{
  pdftitle={Pédagogie des TP IA numérique. Annexe A : choix de conception (support HTML centralisé, formats de fichiers, scénarisation, défi et cybersécurité de l'IA)},
  pdfauthor={Selina Cheggour},
  pdfsubject={Annexe à la méthode « Pédagogie des TP IA numérique ». Présentée au conseil de perfectionnement de l'Icam Toulouse, 14-15 octobre 2026. © 2026 Selina Cheggour, tous droits réservés.},
  pdfkeywords={pédagogie, IA, fichier HTML unique, formats de fichiers, scénarisation, cybersécurité, injection indirecte, tous droits réservés}
}

\color{encre}
\titleformat{\section}{\fontfamily{Montserrat-TLF}\selectfont\Large\bfseries\color{encre}}{\color{icamorange}A.\thesection}{0.6em}{}[\vspace{-0.6em}{\color{icamorange}\rule{18mm}{1.2pt}}]
\titleformat{\subsection}{\fontfamily{Montserrat-TLF}\selectfont\large\bfseries\color{icamdark}}{A.\thesubsection}{0.6em}{}
\titlespacing*{\section}{0pt}{1.4em}{0.9em}

\pagestyle{fancy}
\fancyhf{}
\renewcommand{\headrulewidth}{0pt}
\fancyhead[L]{\scriptsize\color{icamgray}Pédagogie des TP IA numérique · Annexe A}
\fancyhead[R]{\scriptsize\color{icamgray}\auteure}
\fancyfoot[C]{\scriptsize\color{icamgray}\droits\\\presente\quad·\quad page \thepage/\pageref{LastPage}}
\fancypagestyle{plain}{\fancyhf{}\fancyfoot[C]{\scriptsize\color{icamgray}\droits\\\presente\quad·\quad page \thepage/\pageref{LastPage}}}

\newtcolorbox{cadre}[1][]{enhanced,colback=fondorange,colframe=fondorange,boxrule=0pt,arc=0pt,
  left=4mm,right=4mm,top=2.5mm,bottom=2.5mm,borderline west={2pt}{0pt}{icamorange},#1}
\newtcolorbox{cadregris}[1][]{enhanced,colback=fondgris,colframe=fondgris,boxrule=0pt,arc=0pt,
  left=4mm,right=4mm,top=2.5mm,bottom=2.5mm,#1}

\setlist{itemsep=2pt,topsep=4pt}
\setlength{\parindent}{0pt}
\setlength{\parskip}{0.55em}

\begin{document}
\thispagestyle{plain}

{\scriptsize\bfseries\color{icamred}PÉDAGOGIE DES TP IA NUMÉRIQUE · ANNEXE A · 15 OCTOBRE 2026}\par\vspace{3mm}
{\fontfamily{Montserrat-TLF}\selectfont\fontsize{22}{27}\selectfont\bfseries Choix de conception\par}
\vspace{1mm}
{\large\color{icamgray}Support HTML centralisé, formats de fichiers, scénarisation,\\ défi et cybersécurité de l'IA\par}
\vspace{2mm}
{\color{icamgray}\auteure{} · Icam Strasbourg-Europe}\par
\vspace{4mm}
\begin{cadre}
\textbf{\presente.} Cette annexe complète la note de méthode « Pédagogie des TP IA numérique :
étalonner l'instrument, de l'atelier guidé au mini-projet » (DOI 10.5281/zenodo.23390568). La note
décrit les phases et le contenu ; l'annexe décrit les choix de conception qui les portent : le
support, les formats de fichiers et l'esprit du dispositif. L'auteure en revendique la paternité.
\end{cadre}

% ---------------------------------------------------------------
\section{Le support centralisé : l'atelier en un fichier HTML unique}

L'atelier guidé est livré comme \textbf{un seul fichier HTML, autonome et centralisé}. Tout ce
dont l'étudiant a besoin pendant la séance est contenu dans ce fichier : consignes, cas
d'entreprise, documents, données, images, outils d'aide et pièges pédagogiques. Il n'y a rien à
déposer à côté et aucun lien interne qui puisse casser. Le fichier se dépose tel quel sur une
plateforme d'enseignement (Moodle) ou s'ouvre localement dans un navigateur.

\begin{center}
\renewcommand{\arraystretch}{1.4}
\small
\begin{tabularx}{\linewidth}{>{\bfseries\raggedright\arraybackslash}p{4.2cm} >{\raggedright\arraybackslash}X}
\toprule
\textbf{Élément intégré} & \textbf{Rôle dans la méthode}\\
\midrule
Parcours en manches & Cadrage, cinq manches (un outil, un geste d'ingénieur) et une manche de passage au mini-projet ; une manche à la fois, avec « ce qu'il fallait en retenir » à la fin de chacune.\\
Progression mémorisée & La page retient localement, dans le navigateur, jusqu'où l'étudiant est arrivé.\\
Document de référence intégré & Le PDF de la note interne est embarqué dans la page (encodé, reconstruit en fichier à la volée) : visionneuse, téléchargement pour l'outil documentaire, ouverture dans un onglet.\\
Données intégrées & Le relevé de production (CSV) est embarqué : téléchargement ou copie en un clic.\\
Images intégrées & Le logo et le graphique de rebut sont embarqués ; le graphique porte un piège en micro-caractères.\\
Assembleur de prompt & L'étudiant remplit rôle, contexte, contraintes et format ; la page assemble le prompt et le copie.\\
Prompts prêts à copier & Chaque prompt de l'atelier a son bouton « Copier ».\\
Piège invisible & Un paragraphe en blanc sur blanc, récupéré par tout copier-coller de la page vers une IA.\\
\bottomrule
\end{tabularx}
\end{center}

L'intégration du PDF est réalisée par le script \texttt{integrer\_document.py}, écrit par
l'auteure : il encode le document, l'insère avec sa visionneuse et ses boutons, et remplace une
intégration précédente (opération idempotente). Il est déposé avec cette annexe.

% ---------------------------------------------------------------
\section{Les formats de fichiers, choisis pour ce qu'ils enseignent}

Chaque format a été choisi pour une raison pédagogique, et non par commodité. Le format fait
partie de l'exercice : il détermine ce que l'étudiant voit, ce que l'IA lit, et où un piège peut
se loger.

\begin{center}
\renewcommand{\arraystretch}{1.4}
\small
\begin{tabularx}{\linewidth}{>{\bfseries\raggedright\arraybackslash}p{2.6cm} >{\raggedright\arraybackslash}p{3.6cm} >{\raggedright\arraybackslash}X}
\toprule
\textbf{Format} & \textbf{Pour quel support} & \textbf{Ce que ce choix enseigne}\\
\midrule
HTML unique & L'atelier guidé & Un support interactif et autonome ; ce que l'écran montre n'est pas tout ce que la page contient (texte invisible).\\
PDF & L'énoncé, le dossier LC-4, le mémo, la note de référence & Le format « officiel » de l'entreprise, qu'on croit figé ; il transporte pourtant des métadonnées (titre, sujet, mots-clés) lues par les outils et absentes des pages : c'est le piège du mini-projet. Énoncé et dossier sont deux PDF séparés, pour inciter à joindre le dossier à l'IA.\\
PNG & Le graphique de rebut & Une image dont le texte n'est pas sélectionnable : seule la vision de l'IA la lit, y compris des micro-caractères que l'œil ne voit pas. On enseigne de transmettre le fichier d'origine plutôt qu'une capture d'écran.\\
CSV & Le relevé de production & Des données brutes à importer dans un tableur ; une colonne non nommée dans la question (le statut des lignes) suffit à fausser un montant calculé par l'IA.\\
\LaTeX{} & Les sources des documents & Une charte maîtrisée, des documents reproductibles, et des métadonnées choisies au caractère près.\\
\bottomrule
\end{tabularx}
\end{center}

% ---------------------------------------------------------------
\section{L'esprit du dispositif}

\subsection{Scénarisé}
L'étudiant ne fait pas des exercices : il joue un rôle dans une histoire continue. Ingénieur
junior au service méthodes d'une PME fictive (« MécaIcam »), il vit son premier jour : un
dossier à préparer pour un comité qui se réunit vendredi, chaque manche faisant avancer ce
dossier d'un cran. Le mini-projet prolonge l'histoire : une ligne à l'arrêt, trois cents palettes
en attente, une livraison client, une pression commerciale réelle et une réponse signée attendue
avant la fin de la séance. Le récit donne un enjeu à chaque vérification.

\subsection{Challengeant}
\begin{itemize}
\item Les pièges ne sont jamais annoncés ; seuls des réflexes de lecture sont donnés.
\item La difficulté monte : un vecteur de piège jamais vu en atelier attend dans le mini-projet.
\item La décision est binaire et le piège la fait basculer : se tromper a une conséquence visible.
\item L'outil est mis en défaut devant l'étudiant : norme inventée, rétractation au simple
« Tu es sûr ? », montant faux d'aplomb. Un accord obtenu en insistant ne vaut rien.
\item La réussite ne se mesure pas à la réponse de l'IA mais à la capacité de dire d'où vient
chaque affirmation, et ce qu'on n'a pas pu établir.
\end{itemize}

\subsection{Cybersécurité de l'IA}
Le dispositif est une initiation à la sécurité des systèmes d'IA, sans aucune compétence
technique requise :
\begin{itemize}
\item \textbf{l'injection indirecte de prompt}, première vulnérabilité des applications fondées
sur un modèle de langage, vécue de l'intérieur sur trois vecteurs (texte invisible, image,
métadonnées) ;
\item \textbf{l'empoisonnement par l'information} : le piège le plus efficace n'ordonne rien, il
ajoute une donnée plausible au dossier ;
\item \textbf{les défaillances du modèle} : hallucination de références, complaisance envers
l'utilisateur, réponse juste à une question mal posée ;
\item \textbf{l'hygiène des données} : ne jamais confier à une IA un document réel confidentiel
ni des données personnelles ;
\item \textbf{la posture du défenseur} : inspecter ce qu'on transmet à l'IA (tout sélectionner,
lire à voix haute, ouvrir les propriétés du fichier, regarder l'image d'origine) et tracer ce
qu'elle renvoie.
\end{itemize}

% ---------------------------------------------------------------
\section{Éléments originaux revendiqués}

L'auteure revendique, en complément de la note de méthode, dans leur formulation et leur
combinaison :
\begin{enumerate}
\item la livraison d'un atelier pédagogique sur l'IA sous la forme d'un fichier HTML unique,
autonome et centralisé, intégrant consignes, documents, données, images, outils et pièges ;
\item l'intégration d'un document de référence PDF dans cette page, consultable et exportable
vers un outil d'IA documentaire ;
\item le choix raisonné des formats (HTML, PDF, PNG, CSV) comme vecteurs pédagogiques, chacun
porteur d'un enseignement et, le cas échéant, d'un piège propre à ce format ;
\item la scénarisation continue (un rôle, une entreprise fictive, un comité, une ligne à
l'arrêt) reliant l'atelier et le mini-projet ;
\item la conception challengeante : pièges non annoncés, difficulté croissante, décision qui
bascule, évaluation par la traçabilité ;
\item l'approche « cybersécurité de l'IA » pour débutants, fondée sur l'injection indirecte et
l'empoisonnement par l'information ;
\item l'assembleur de prompt structuré, le parcours en manches à progression mémorisée, le
script \texttt{integrer\_document.py} et la page \texttt{TP5-IA.html}.
\end{enumerate}

\vfill
\begin{cadregris}
\small Annexe établie par \auteure{} le 15 octobre 2026. \presente. \droits.
\end{cadregris}
\end{document}
